%----------------------------------------------------------------
%
%  File    :  survey-introduction.tex
%
%  Author  :   Lukas Auer, David Heidinger, Nina Tschikof and Christina Vogel, TU Graz, Austria
% 
%  Created :  06 May 2026
% 
%  Changed :  06 May 2026
% 
%----------------------------------------------------------------


\chapter{Introduction}

\label{chap:Intro}

\section{Introduction}

Data visualisations are widely used to communicate complex information in an intuitive and efficient way. They support decision-making across domains such as science, journalism, business or politics. However, the same properties that make visualisations powerful also make them susceptible to misuse.

Dishonest chart techniques refer to design choices or data manipulations that distort the viewer's perception of the underlying data. These techniques may occur either intentionally or unintentionally as a result of poor design choice and can significantly influence how patterns, trends and relationships are interpreted.

This paper provides a survey of common dishonest chart techniques and proposes a taxonomy of different techniques based on their characteristics. It also discusses their impact on data interpretation. Additionally, related work and strategies for identifying and avoiding misleading visualisations are reviewed.

This paper aims to provide a taxonomy of common dishonest chart techniques and to analyse their underlying mechanisms. By categorising these techniques it becomes easier to identify patterns, understand how different forms of distortion operate and what to look out for. By increasing awareness of dishonest practices both creators and consumers of visualisations can contribute to more accurate and transparent communication of information.



